\scanpage{116}
% ms. line 01
Soit $(r_{\unsure{\nu|v}})$ suite croissante de nombres \deleted{\unsure{positifs}} $>0$, $\to\infty$
% ms. line 02
soit $n(r)$ le nombre des indices $\nu$ tels que
% ms. line 03
$r_\nu\leq r$, on a, pour $0\leq a\leq b<+\infty$, et $\alpha>0$
% ms. line 04
\[ \boxed{\displaystyle \sum_{\substack{\nu\\(a<r_\nu<b)}}\frac1{r_\nu^\alpha}=\left(\frac{n(b-0)}{b^\alpha}-\frac{n(a+0)}{a^\alpha}\right)+\alpha\int_a^b\frac{n(r)\,dr}{r^{\alpha+1}}}\tag*{(1)} \]
% ms. line 05
(car le premier membre est égal à $\int_{\unsure{a_0|r_0|a}}^b r^{-\alpha}\,dn(r)$,
% ms. line 06
et on intègre par parties)

% ms. line 07
\uline{Proposition} \quad $\displaystyle \sum_\nu\frac1{r_\nu^\alpha}<+\infty\ \Longleftrightarrow\ \int_0^\infty\frac{n(r)\,dr}{r^{\alpha+1}}\unsure{\leq|<}+\infty$
% ms. line 08
alors $\frac{n(r)}{r^\alpha}\to0$ pour $r\to\infty$, et on a
% ms. line 09
\[ \boxed{\displaystyle \sum_{\substack{\nu\\(r_\nu>a)}}\frac1{r_\nu}=-\frac{n(a+0)}{a^\alpha}+\alpha\int_a^\infty\frac{n(r)\,dr}{r^{\alpha+1}}}\tag*{\unsure{(1)|(1\textquotesingle)}} \]
% Vertical stroke in the left margin alongside the proposition.

% ms. line 10
\deleted{$\leftarrow$ \unsure{si} $\sum_\nu\frac1{r_\nu^\alpha}<+\infty$, on a $\frac{n(r)}{r^\alpha}\leq M$ pour \unsure{tout} $r$
% ms. line 11
où $M=$ constante}
% ms. line 12
En effet, (1) s'écrit, pour $a=0$, \deleted{\unsure{comme}}
% ms. line 13
\[ \sum_{r_\nu<b}\frac1{r_\nu^{\unsure{\alpha|\nu}}}=\frac{n(b-0)}{b^\alpha}+\alpha\int_0^b\frac{n(r)\,dr}{r^{\alpha+1}}\tag*{\textcircled{\scriptsize 1 bis}} \]
% ms. line 14
d'où suit que si \deleted{\unsure{pour un} $b$} $\sum_\nu\frac1{r_\nu^\alpha}$ est
% ms. line 15
convergente, il en est de $=$ de $\int_0^\infty\frac{n(r)\,dr}{r^{\alpha+1}}$
% ms. line 16
Réciproquement, si $\int_{\unsure{0}}^\infty\frac{n(r)\,dr}{r^{\alpha+1}}<+\infty$ alors il
% ms. line 17
existe $M>0$ et une suite croiss. de $b_i\to\infty$
% ms. line 18
de nbs tels que $\frac{n(b_i)}{b_i^\alpha}\leq M$ ; \unsure{p.} on a donc
% ms. line 19
$\displaystyle \sum_{r_\nu<b_i}\frac1{r_\nu^\alpha}\leq M+\alpha\int_0^\infty\frac{n(r)\,dr}{r^{\alpha+1}}$, d'où $\displaystyle \tikz[baseline=(a.base)]{\node[inner sep=1pt] (a) {$\displaystyle\sum_{\unsure{i|\nu}}\frac1{r_i^\alpha}$};\draw ([yshift=3pt]a.south west)--(a.south west)--(a.south east)--([yshift=3pt]a.south east);}<+\infty$.
% archive sheet number: 117
